\chapter{Problem Definition and Mathematical Formulation}
\label{chap:problem}

% ---------------------------------------------------------------
\section{Chapter Overview and Notation}
% ---------------------------------------------------------------
\begin{table}[ht]
\centering
\small
\begin{tabular}{l>{\raggedright\arraybackslash}p{9.8cm}}
\toprule
\textbf{Symbol}                        & \textbf{Description} \\
\midrule
$t$                                    & Time index (hourly intervals) \\
$\Delta t = 1$\,h                      & Time-step duration \\
$P_t^{\PV}$                            & Solar PV generation (kW) \\
$P_t^{\Load}$                          & Telecom load demand (kW) \\
$G_t \in \{0,1\}$                      & Grid availability indicator \\
$\SoC_t \in [0,1]$                     & Battery state-of-charge (normalised) \\
$\Inv_t$                               & On-site diesel inventory (litres) \\
$\Pipe_t$                              & Total diesel quantity in transit (litres); at most one
                                         pending order at any time \\
$P_t^{\Bat}$                           & Battery power (positive = charge, negative = discharge) (kW) \\
$P_{\max}^{\ch},\; P_{\max}^{\dis}$    & Battery charge/discharge power limits (kW) \\
$\DG_t \in \{0,1\}$                    & Diesel generator ON/OFF decision (1 = ON at fixed power) \\
$P_{\mathrm{rated}}^{\DG}$             & DG fixed operating power when on (kW) \\
$P_t^{\DG}$                            & DG electrical output $= \DG_t \cdot P_{\mathrm{rated}}^{\DG}$ \\
$Q_t$                                  & Diesel order quantity (litres) \\
$r_{\DG}$                              & DG fuel consumption rate (litres/kWh) \\
$C^{\Tank}$                            & Tank capacity (litres) \\
$E_{\Bat}$                             & Battery energy capacity (kWh) \\
$\eta_{\ch},\; \eta_{\dis}$            & Charge/discharge efficiency \\
$\hour_t$                              & Hour-of-day encoded as $(\sin(2\pi h/24),\cos(2\pi h/24))$ \\
$L_t$                                  & Stochastic lead time (timesteps) for a placed diesel order \\
$\tau_t$                               & Hours elapsed since the last order was placed (reset to 0 on each new order) \\
$R_t$                                  & Estimated hours remaining until delivery; identically zero except in the
                                         separate lognormal-forecast extension not used in the main experiments \\
$A_t$                                  & Diesel arrival quantity at time $t$ (litres) \\
$\gamma$                               & Discount factor \\
\midrule
\multicolumn{2}{l}{\textit{Objective function coefficients}} \\
\midrule
$P_t^{\Grid}$                          & Grid power drawn at time $t$ (kW) \\
$P_t^{\Grid,\max}$                     & Maximum grid connection capacity (kW) \\
$c_{\Grid}$                            & Grid electricity cost coefficient (Rs./kWh) \\
$c_{\mathrm{diesel}}$                  & Diesel fuel cost (Rs./litre) \\
$c_{\order}$                           & Fixed administrative cost per order placed (Rs.) \\
$c_{\hold}$                            & Inventory holding cost coefficient (Rs./litre/h) \\
$c_{\degrad}$                          & Battery degradation cost coefficient (Rs./kWh throughput) \\
$c_{\COtwo}$                           & Carbon cost coefficient (Rs./kg\,CO$_2$e) \\
$\EF_{\mathrm{diesel}}$                & Diesel emission factor ($= 2.68$\,kg\,CO$_2$e/litre) \\
$\lambda_{\short}$                     & Penalty coefficient for fuel-shortage events \\
\bottomrule
\end{tabular}
\caption{Complete notation used in the mathematical formulation.}
\label{tab:notation}
\end{table}

As established in Chapter~\ref{chap:literature} and summarised in
Table~\ref{tab:litgap}, to the best of the author's knowledge no prior work
jointly models diesel fuel as a finite stochastic inventory while
simultaneously optimising operational energy dispatch in telecom towers.
In rural deployments, diesel logistics---including pilferage, stockouts, and
delivery disruptions---contribute significantly to tower outages
\citep{datoms2024,alton2025,gsma2025b,trai2024}.

This chapter addresses this gap by presenting a unified Constrained Markov
Decision Process (CMDP) formulation that treats battery state-of-charge and
diesel inventory as jointly controlled stochastic stocks under uncertainty.
All decisions are taken at hourly resolution ($\Delta t = 1$\,h), consistent
with the ITU/Zindi dataset used in this work (described in
Chapter~\ref{chap:datasets-setup}).


% ---------------------------------------------------------------
\section{System Model}
% ---------------------------------------------------------------

The overall hybrid telecom tower and diesel-logistics system is illustrated in
Figure~\ref{fig:system_diagram}.

\begin{figure}[htbp]
\centering
\resizebox{0.9\textwidth}{!}{%
\begin{tikzpicture}[
    >=Stealth,
    node distance=2.6cm and 3.2cm,
    box/.style={draw, rounded corners=4pt, minimum height=1.1cm,
                minimum width=2.8cm, align=center, font=\small},
    agent/.style={draw, fill=gray!12, rounded corners=6pt,
                  minimum width=4.5cm, minimum height=1.4cm,
                  align=center, font=\bfseries\small},
    arrow/.style={->, thick},
    obs/.style={->, dashed, thick, blue!70!black},
    act/.style={->, dashed, thick, red!70!black}
]

% --- Energy components (top row) ---
\node[box] (pv)   {Solar PV};
\node[box, right=of pv]   (grid) {Grid Supply\\(intermittent)};
\node[box, right=of grid] (bat)  {Battery Storage\\(LFP)};
\node[box, right=of bat]  (dg)   {Diesel Generator\\(fixed $P_{\mathrm{rated}}^{\DG}$)};
\node[box, right=of dg]   (load) {Telecom Load\\(BTS + ancillary)};

% --- RL agent ---
\node[agent, below=2.8cm of bat] (agent) {RL Agent};

% --- Diesel logistics (bottom row) ---
\node[box, below=5.0cm of dg] (tank) {Diesel Tank\\($\Inv_t$ litres)};
\node[box, left=of tank]       (pipe) {Pipeline\\($\Pipe_t$)};
\node[box, left=of pipe]       (sup)  {Fuel Supplier};

% --- Physical flows ---
\draw[arrow] (pv)   -- (bat);
\draw[arrow] (grid) -- (bat);
\draw[arrow] (bat)  -- (load);
\draw[arrow] (dg)   -- (load);
\draw[arrow] (sup)  -- node[above,font=\footnotesize] {$Q_t$ (order)} (pipe);
\draw[arrow] (pipe) -- node[above,font=\footnotesize] {$A_t$ (arrival)} (tank);
\draw[arrow] (tank) -- node[right,font=\footnotesize] {fuel} (dg);

% --- Observations ---
\draw[obs] (pv.south)   -- ++(0,-1.2) -| (agent.north west);
\draw[obs] (grid.south) |- (agent.north);
\draw[obs] (bat.south)  -- (agent.160);
\draw[obs] (load.south) |- (agent.north east);
\draw[obs] (pipe.north) -- ++(0,1.8) -| (agent.south west);

\node[font=\footnotesize, above=7mm of agent, align=left]{
  \textbf{State $s_t$:} PV output $P_t^{\PV}$, load $P_t^{\Load}$,\\
  grid availability $G_t$, battery SoC,\\
  diesel inventory $\Inv_t$, pipeline $\Pipe_t$,\\
  time-of-day encoding $\hour_t$
};

% --- Actions ---
\draw[act] (agent.60)  -- node[right,font=\footnotesize] {$\DG_t \in \{0,1\}$}    (dg.south);
\draw[act] (agent.240) -- node[left,font=\footnotesize]  {$Q_t$}                   (sup.north);

\end{tikzpicture}%
}
\caption{System architecture of the hybrid telecom tower under DG ON/OFF
modelling. The RL agent observes energy states (PV, load, grid, battery SoC)
and diesel logistics states (tank inventory $\Inv_t$ and in-transit pipeline
$\Pipe_t$), and controls DG ON/OFF via
$\DG_t \in \{0,1\}$ and diesel order quantity $Q_t$. Battery dispatch
$P_t^{\Bat}$ is determined by an internal environment rule rather than a
direct RL action.}
\label{fig:system_diagram}
\end{figure}

The system comprises rooftop solar PV, a lithium iron phosphate (LFP) battery
bank, intermittent grid supply, a diesel generator operating at fixed power
$P_{\mathrm{rated}}^{\DG}$ when switched on, and a finite diesel tank
replenished by external truck deliveries subject to stochastic lead times.
$G_t \in \{0,1\}$ is a binary grid availability indicator; when $G_t = 0$ the
grid is unavailable due to outage and no power can be drawn from it.
Decisions are taken every hour ($\Delta t = 1$\,h). The key modelling
assumptions are: perfect state observability within each timestep;
non-cancellable diesel orders once placed; and a minimum DG up-time of 2\,h
(2~steps at hourly resolution) whenever the generator is switched on.
The present study focuses on single-tower energy management using the
ITU/Zindi dataset (10 sites, hourly resolution, 60 days each); multi-tower
coordination, site-specific hardware variation, and long-horizon logistics
network optimisation are outside the scope of the current work.

At each timestep, the power balance constraint must be satisfied:
\begin{equation}
P_t^{\Load} = P_t^{\PV} - P_t^{\Bat} + P_t^{\Grid} + \DG_t\,P_{\mathrm{rated}}^{\DG},
\label{eq:balance}
\end{equation}
where $P_t^{\Grid} = G_t \cdot P_t^{\Grid,\max}$ is the grid power drawn, bounded
by grid availability and the local connection capacity. The negative sign on
$P_t^{\Bat}$ reflects the convention of Table~\ref{tab:notation}: a positive
value denotes charging (power absorbed from the bus), so a discharging battery
contributes positively to the supply side via the negative sign.

Although diesel inventory decisions operate on multi-day timescales, they
directly constrain hourly dispatch feasibility. During grid outages or
low-solar periods, the generator can only operate if sufficient fuel remains
on site. A controller that optimises hourly energy dispatch without accounting
for inventory state may therefore select locally optimal actions that lead to
infeasible states in subsequent hours. Jointly modelling inventory and dispatch
decisions allows the agent to anticipate future fuel constraints and maintain
service reliability under uncertain delivery schedules.

The formulation therefore differs from conventional telecom energy-management
models in two respects. First, diesel fuel is represented as a stochastic
inventory subject to delivery delays rather than as an unconstrained fuel
source. Second, dispatch and inventory decisions are optimised jointly within
a single sequential decision framework, allowing the controller to anticipate
future fuel availability when making present operational decisions.

% ---------------------------------------------------------------
\section{Motivation for Reinforcement Learning}
\label{sec:rl_motivation}
% ---------------------------------------------------------------

Although the system dynamics can be written analytically, the resulting
control problem becomes a stochastic mixed-integer dynamic program once
diesel ordering, generator switching, and dispatch decisions are
coupled across time, requiring accurate probabilistic models of solar
generation, load, grid outages, and delivery delay that are difficult
to specify in closed form and vary across sites and seasons.
Reinforcement learning avoids solving this optimisation problem from
scratch at every step by learning a policy directly from interaction,
and the trained policy generalises across sites and scenarios without
re-solving. The algorithmic choices this motivates (PPO, action
masking, and the unified-vs-decomposed architecture question) are
developed in Chapter~\ref{chap:methodology}.

% ---------------------------------------------------------------
\section{State Space}
\label{sec:state}
% ---------------------------------------------------------------

The state vector observed by the agent at each hourly step is
\begin{equation}
s_t = \bigl[\,
  \SoC_t,\;
  \Inv_t,\;
  \Pipe_t,\;
  P_t^{\PV},\;
  P_t^{\Load},\;
  G_t,\;
  \hour_t,\;
  \tau_t,\;
  R_t
\,\bigr],
\label{eq:state}
\end{equation}
where $\hour_t = (\sin(2\pi h/24),\,\cos(2\pi h/24))$ encodes the hour of
day as a two-dimensional cyclic feature, $\tau_t$ is the elapsed time
since the last order was placed, and $R_t$ is an estimated
delivery-remaining-time signal that is identically zero throughout the
main experiments (Chapter~\ref{chap:results}) and becomes informative
only in a separate lognormal-forecast extension outside the primary
scope of this thesis.

The observation consists of nine logical components, corresponding to
eleven scalar inputs after expansion of the cyclic hour-of-day encoding.
Table~\ref{tab:state_sources} documents the data provenance of
each component, a distinction the panel raised explicitly in Phase-I review.

\begin{table}[ht]
\centering
\small
\begin{tabular}{lll}
\toprule
\textbf{Variable} & \textbf{Source} & \textbf{Notes} \\
\midrule
$P_t^{\PV}$   & ITU/Zindi dataset  & Direct column \\
$P_t^{\Load}$ & ITU/Zindi dataset  & Direct column \\
$G_t$         & ITU/Zindi dataset  & Direct column \\
$\SoC_t$      & Simulator state    & Evolved via Eq.~\eqref{eq:soc} from real initial values \\
$\Inv_t$      & Simulator state    & Derived from ITU fuel-level column and tank capacity \\
$\Pipe_t$     & Synthetic          & Log-normal lead-time overlay; sole synthetic component \\
$\hour_t$     & Derived            & Computed from ITU timestamp \\
$\tau_t$      & Simulator state    & Counter reset on each new order; always active \\
$R_t$         & Synthetic (inactive) & Present in the observation vector but fixed at zero for
                                       every experiment reported in this thesis \\
\bottomrule
\end{tabular}
\caption{Data provenance of each state variable. $\Pipe_t$ and $R_t$ are
introduced synthetically; $R_t$ additionally never takes a non-zero
value anywhere in the reported experiments. The remaining components are
either obtained directly from the ITU/Zindi dataset or evolved from
dataset-grounded initial conditions.}
\label{tab:state_sources}
\end{table}

The pipeline quantity is capped at one pending order at any time: if a
delivery is already in transit, the feasible action set restricts $Q_t = 0$
until the pending order arrives (Section~\ref{sec:constraints}).

% ---------------------------------------------------------------
\section{Action Space}
\label{sec:action}
% ---------------------------------------------------------------

At each hourly step the agent selects the control actions
\begin{equation}
a_t = \bigl[\,
  \DG_t,\;
  Q_t
\,\bigr],
\label{eq:action}
\end{equation}
with
\[
\DG_t \in \{0,1\}, \qquad
Q_t \in \{0,\,500,\,1000,\,1500,\,2000\}\;\text{litres.}
\]

Here, $\DG_t$ controls generator commitment and $Q_t$ controls diesel
replenishment. Battery dispatch $P_t^{\Bat}$ is determined by an
internal environment rule that prioritises renewable and battery energy
before diesel generation, rather than being a direct RL action.

The five discrete order quantities reflect standard tanker delivery sizes used
by Indian telecom operators, covering the practical range of replenishment
volumes. The action space therefore consists of two discrete control
decisions: generator commitment and diesel order quantity. Infeasible
actions are excluded from the feasible set $\mathcal{C}(s_t)$ defined in
Section~\ref{sec:cmdp}; the implementation of this feasible-set
enforcement is described in Chapter~\ref{chap:methodology}.

% ---------------------------------------------------------------
\section{System Dynamics and Uncertainty}
\label{sec:dynamics}
% ---------------------------------------------------------------

\paragraph{Battery dynamics.}
Battery state-of-charge evolves according to
\begin{equation}
\SoC_{t+1} = \SoC_t
  + \frac{\eta_{\ch}\,\max(P_t^{\Bat},0)\,\Delta t
          - \max(-P_t^{\Bat},0)\,\Delta t\,/\,\eta_{\dis}}
         {E_{\Bat}}.
\label{eq:soc}
\end{equation}

\paragraph{Diesel inventory dynamics.}
On-site fuel inventory evolves as
\begin{equation}
\Inv_{t+1} = \min\!\bigl(C^{\Tank},\;
  \Inv_t - r_{\DG}\,P_t^{\DG}\,\Delta t + A_t\bigr),
\label{eq:inv}
\end{equation}
where the arrival term is
\begin{equation}
A_t =
\begin{cases}
  Q_{t-L_t}, & \text{if a pending order's lead time expires at step } t, \\
  0,          & \text{otherwise,}
\end{cases}
\label{eq:arrival}
\end{equation}

\paragraph{Pipeline dynamics.}
The pipeline state $\Pipe_t$ represents the quantity of diesel currently in
transit from a previously placed order. When the sampled lead time expires,
the corresponding order quantity $Q_{t-L_t}$ is transferred to on-site
inventory via the arrival term $A_t$ (Equation~\eqref{eq:arrival}); until
that point $\Pipe_t = Q_{t-L_t}$ and no further orders may be placed. In
the absence of a pending order, $\Pipe_t = 0$.

\paragraph{Stochastic lead time.}
Lead time in days is drawn from a log-normal distribution at the moment
an order is placed:
\begin{equation}
L_t^{\text{days}} \;\sim\; \LogNormal(\mu,\,\sigma=0.6),
\label{eq:leadtime}
\end{equation}
with $\mu = 2.5$\,days under the \emph{normal logistics} scenario and
$\mu = 5.0$\,days under the \emph{delayed logistics} scenario. Both
scenarios are evaluated in Chapter~\ref{chap:results} to test robustness
under supply-chain disruption. The lead time in hourly timesteps is
\begin{equation}
L_t = \lceil 24\,L_t^{\text{days}} \rceil,
\label{eq:leadtime_steps}
\end{equation}
consistent with $\Delta t = 1$\,h.

The log-normal distribution for $L_t^{\text{days}}$ is standard in remote
diesel logistics modelling \citep{srinivasan2013,shen2011}. Because actual
delivery records are unavailable in the ITU/Zindi dataset, $\Pipe_t$ is the
only state component generated synthetically; all other dynamics are driven
directly by real sensor measurements.

% ---------------------------------------------------------------
\section{Constraints}
\label{sec:constraints}
% ---------------------------------------------------------------

Three categories of hard constraint are enforced with zero tolerance in
closed-loop operation.

The \emph{battery constraints} require that state-of-charge remains within
the safe operating window at all times, i.e.\ $0.20 \leq \SoC_{t+1} \leq 1.00$
(corresponding to 80\% maximum depth-of-discharge for LFP cells), and that
the charging and discharging power satisfies
$-P_{\max}^{\dis} \leq P_t^{\Bat} \leq P_{\max}^{\ch}$. The feasible battery
state transitions are obtained analytically from Equation~\eqref{eq:soc}
and constrain the internal battery-dispatch rule at each step.

The \emph{generator constraint} requires that the DG may only be switched on
if sufficient fuel remains to sustain at least 2\,h of continuous operation
at rated power, i.e.\ $\Inv_t \geq r_{\DG}\,P_{\mathrm{rated}}^{\DG}\cdot 2$.
If this condition is not met, $\DG_t = 1$ is infeasible in that state.
Additionally, once the generator is switched on it must remain on for a
minimum of 2\,h (2~steps), reflecting realistic start-up costs and equipment
limitations.

The \emph{ordering constraint} prevents a new diesel order from being placed
while a delivery is already in transit ($\Pipe_t > 0$), making any $Q_t > 0$
infeasible in that state, and caps the maximum order quantity at remaining
tank headroom, i.e.\ $Q_t \leq C^{\Tank} - \Inv_t$. Together, the three
constraint categories define the feasible action set $\mathcal{C}(s_t)$
incorporated into the CMDP tuple of Section~\ref{sec:cmdp}. Their
implementation in the training algorithm is described in
Chapter~\ref{chap:methodology}.

% ---------------------------------------------------------------
\section{Objective Function}
\label{sec:objective}
% ---------------------------------------------------------------

The agent minimises the expected total discounted cost over an episode:
\begin{align}
J(\pi) = \mathbb{E}\Biggl[\,\sum_{t=0}^{T} \gamma^t \Bigl(
    &\;c_{\Grid}\,P_t^{\Grid}
     + c_{\mathrm{diesel}}\,r_{\DG}\,P_t^{\DG}\,\Delta t
     + c_{\order}\,\mathbf{1}(Q_t > 0)
     \notag \\[4pt]
    &+ c_{\hold}\,\Inv_t
     + c_{\degrad}\,|P_t^{\Bat}|
     \notag \\[4pt]
    &+ c_{\COtwo}\,\EF_{\mathrm{diesel}}\,r_{\DG}\,P_t^{\DG}\,\Delta t
     \notag \\[4pt]
    &+ \lambda_{\short}\,\max\!\bigl(0,\;
       r_{\DG}\,P_t^{\DG}\,\Delta t - \Inv_t\bigr)
\Bigr)\Biggr],
\label{eq:objective}
\end{align}
where $\gamma = 0.995$ gives an effective planning horizon of approximately
200\,h ($\approx 8$\,days), sufficient to cover realistic lead times.

The cost terms are: $c_{\Grid}$, the grid electricity cost coefficient;
$c_{\mathrm{diesel}}$, the diesel fuel cost in Rs./litre; $c_{\order}$, a
fixed administrative cost per order placed; $c_{\hold}$, an inventory
holding cost coefficient; $c_{\degrad}$, a battery degradation cost
proportional to total charge throughput; $c_{\COtwo}$, a carbon cost
coefficient (Rs.\ per kg CO$_2$e); $\EF_{\mathrm{diesel}} = 2.68$\,kg\,CO$_2$e/litre,
the diesel emission factor; and $\lambda_{\short}$, a large penalty
coefficient applied whenever the generator would consume more fuel than is
available on-site. All cost coefficients are calibrated from published Indian telecom
operator data \cite{IOCL_Prices2024, SERC_Tariff2024, Operator_Diesel2023}.

The soft stockout penalty ($\lambda_{\short}$ term) discourages the policy
from approaching the fuel-shortage boundary in states where the feasibility
constraint cannot be applied analytically. In the constrained-site evaluation
reported in Chapter~\ref{chap:results}, no hard violations were observed during
policy rollout.

Although the optimisation objective is expressed as discounted cumulative
cost, policy performance is evaluated in Chapter~\ref{chap:results} using
operational metrics including Expected Energy Not Served (EENS), diesel
consumption, network uptime, and a cost proxy, which together reflect the
trade-offs that matter to telecom tower operators.

% ---------------------------------------------------------------
\section{Formal CMDP Definition}
\label{sec:cmdp}
% ---------------------------------------------------------------

The hybrid telecom energy-management problem is modelled as a discounted
episodic Constrained Markov Decision Process defined by the tuple
\[
(\mathcal{S},\,\mathcal{A},\,\mathcal{P},\,c,\,\mathcal{C},\,\gamma),
\]
where $\mathcal{S} \subset \mathbb{R}^{11}$ is the state space defined in
Section~\ref{sec:state}, $\mathcal{A}$ is the action space of
Section~\ref{sec:action}, $\mathcal{P}$ captures the stochastic battery and
inventory dynamics of Section~\ref{sec:dynamics}, $c$ is the cost function
of Equation~\eqref{eq:objective}, and $\gamma = 0.995$. Operational
constraints define a feasible action set $\mathcal{C}(s_t)$ for each state;
how these feasibility requirements are enforced during policy learning is
detailed in Chapter~\ref{chap:methodology}.

Training episodes run for 720 steps (30 days), ensuring that the policy
observes multiple fuel-delivery cycles and representative solar variability
patterns. Evaluation episodes use the held-out 15-day test split of the
ITU/Zindi dataset to assess policy performance on unseen operating conditions.

% ---------------------------------------------------------------
\section{Unified versus Decomposed Architecture}
\label{sec:unified_vs_hier}
% ---------------------------------------------------------------

A natural modelling question is whether the diesel ordering decision
($Q_t$, multi-day timescale) and the energy dispatch decisions (generator
commitment and battery behaviour, hourly timescale) should be optimised
jointly within a single agent or decomposed into separate controllers
operating at their natural timescales. This question was raised explicitly
in the Phase-I reviewer comments and is directly addressed by the
experimental design of this work.

The unified formulation presented in this chapter---in which a single
joint agent controls diesel ordering and generator commitment while battery
behaviour is handled by the environment---is referred to as \textbf{RLInv}
throughout. The alternative \textbf{Track~B} architecture pairs a classical
$(s,S)$ inventory policy for diesel ordering with a separate PPO agent for
dispatch, allowing each sub-problem to be solved at its natural timescale.
The two architectures are evaluated empirically in
Chapter~\ref{chap:results}.

The comparison between unified and decomposed architectures allows the
following hypotheses to be tested.

\textbf{H1 (Coupling hypothesis).} Joint optimisation of diesel ordering and
energy dispatch by a single agent provides measurable reliability or
operational cost benefits relative to decomposed control, in which
replenishment and dispatch decisions are managed separately.


\textbf{H2 (Action masking hypothesis).} Incorporating feasibility-aware
action masking improves reliability under stressed logistics when diesel
replenishment is externally managed, while providing little measurable
benefit when replenishment is learned jointly.

\textbf{H3 (Inventory-state hypothesis).} Including explicit diesel inventory
state in the observation improves ordering decisions, relative to an ablated
agent in which the inventory- and delivery-related observation channels are
removed from the observation vector.

These hypotheses are stated here, at the level of the problem formulation,
because they directly shape the choice of state space, action space, and
ablation variants. The formulation above is intentionally designed so that
the role of inventory state, safety constraints, and joint optimisation can
each be isolated and empirically tested through the comparisons defined above.
Results for all three hypotheses are reported in Chapter~\ref{chap:results},
and conclusions are drawn in Chapter~\ref{chap:conclusion} regardless of the
direction of the findings.

% ---------------------------------------------------------------
\section{Research Objectives and Contributions}
\label{sec:objectives}
% ---------------------------------------------------------------

The formulation presented in this chapter serves three related objectives:
it enables joint dispatch and diesel-ordering control within a single
sequential decision framework; it incorporates stochastic delivery-pipeline
dynamics into the state representation, capturing replenishment uncertainty
that existing energy-management formulations do not model; and it encodes
hard operational feasibility requirements directly in the decision problem,
rather than relying on penalty terms alone.

The primary contribution of this chapter is the formulation of telecom tower
energy management as a constrained sequential decision problem that jointly
models diesel inventory logistics and energy dispatch. The chapter defines the
unified state--action representation, the stochastic delivery-pipeline model,
and the experimental hypothesis structure used to evaluate the formulation.
The empirical evaluation of these hypotheses is presented in
Chapter~\ref{chap:results}.

\paragraph{Stakeholder impact.}
Rural telecom infrastructure connects remote populations to emergency
communications, mobile banking, and digital public services.  India operates
several hundred thousand remote towers, many in weak-grid areas where diesel
logistics account for 30--40\% of operational expenditure~\citep{gsma2025b,trai2024}.
Tower outages in these settings directly interrupt critical services and
disproportionately affect underserved communities.  Improving energy and
inventory management at these sites therefore has direct implications for
service reliability, operating cost, and carbon emissions.  The formulation
presented in this chapter targets precisely these drivers by jointly
optimising dispatch and replenishment decisions under realistic supply-chain
uncertainty.